\chapter{Literature Review}
\label{ch:literature}

This chapter reviews the three bodies of work on which the thesis builds: the economics of player valuation (Section~\ref{sec:lit-valuation}), fuzzy clustering methodology (Section~\ref{sec:lit-fuzzy}), and the application of clustering to football players (Section~\ref{sec:lit-football}). Section~\ref{sec:lit-gap} identifies the gap between them and discusses the alternatives to the approach adopted here, and Section~\ref{sec:lit-questions} states the research questions.

\section{Player valuation and the transfer market}
\label{sec:lit-valuation}

The transfer market has become an industry in its own right. In 2019, 18,080 players moved between clubs in different countries, and clubs paid a total of USD~7.35~billion to recruit them \cite{franceschi2024determinants}. Since the Bosman ruling, players have been treated as depreciable intangible assets when clubs invest in their recruitment, which makes the valuation of players a question of financial as well as sporting importance \cite{franceschi2024determinants}.

\citeA{franceschi2024determinants} provide the most comprehensive synthesis of the empirical literature on this question. Their systematic review covers 29 peer-reviewed studies and 111 model specifications, and it classifies the explanatory variables into six categories: time, labour (for instance contract length), performance, player characteristics, club characteristics, and popularity. Several regularities emerge. Age is a recurrent regressor with a hump-shaped relationship to value: among the specifications in which the terms are significant at the 5\% level, the linear term is positive in 32 of 45 and the squared term is negative in 41 of 45, so that value rises with age up to a turning point and falls afterwards. Goals and assists are positively associated with value in every specification in which they are significant, as is playing time, and playing position is a standard control. The authors also warn that the abundance of on-the-field performance variables creates potential multicollinearity, because absolute performance metrics all depend on playing time. Finally, they argue that valuation research should look beyond the five major European leagues: the Portuguese Liga is named among the leagues that had been studied in only one article, even though Benfica and Porto are, in their assessment, very important actors in player development and in the transfer market.

Most of this literature models either transfer fees or market values published by Transfermarkt, a platform on which a community of users estimates the value of players. \citeA{muller2017beyond} use multilevel regression on 4,217 players from the top five European leagues over six seasons to show that data-driven estimates of market value can be comparably accurate to the crowd's. Their discussion of the crowd's limitations is directly relevant to the data used here. Transfermarkt usually updates a player's value only every six to twelve months, so the estimates do not follow performance closely. In addition, crowd estimates tend to be more accurate for players who are known to a large audience, and the activity of the community is lower in some countries and leagues, which limits their usefulness for scouting players in smaller leagues. \citeA{poli2024statistical} estimate a multiple linear regression on 8,389 transfers and reach an $R^2$ of 84.8\%, drawing on the same families of variables, among them the age of the player, the characteristics of the selling and the buying club, and the length of the remaining contract.

Two features of this literature matter for the present thesis. First, these studies relate the \emph{level} of a player's value or fee to his current characteristics. None of them asks whether the way in which a player's multivariate performance profile \emph{changes} is informative about changes in his value once current performance is controlled for. Second, the evidence is concentrated on the big five leagues. A league such as the Primeira Liga, which supplies many players to larger leagues, has received little attention.

\section{Fuzzy clustering}
\label{sec:lit-fuzzy}

Cluster analysis partitions observations into groups that are internally similar. Hard algorithms such as $k$-means assign every observation to exactly one cluster, which imposes a sharp boundary even where the data show none. Fuzzy clustering relaxes this requirement by assigning each observation a degree of membership in every cluster; memberships lie between zero and one and sum to one over the clusters. The Fuzzy C-Means (FCM) algorithm of \citeA{bezdek1984fcm} minimises a generalised least-squares objective by alternating between updating the cluster prototypes and the memberships. Its behaviour is governed by the fuzziness exponent $m>1$: values close to one yield almost crisp partitions, whereas larger values spread the membership more evenly across clusters. In practice convergence is reached in 10 to 25 iterations, and partition-based indices such as the partition coefficient and the partition entropy are used to judge a solution \cite{bezdek1984fcm}.

Three extensions are relevant here. First, in \emph{fuzzy $C$-medoids} the prototype of each cluster is an actual observation rather than an average. This allows dissimilarities that are not Euclidean, such as those needed for mixed-type data, and it makes the prototypes interpretable: each is a real player \cite{krishnapuram2001lowcomplexity}. For medoid-based methods, low values of $m$ between one and about 1.5 are recommended, because with larger values the objective becomes insensitive to the choice of the medoids \cite{durso2022robust}. Second, \emph{noise clustering} adds an extra cluster that lies at a fixed distance from every observation, so that observations far from all genuine clusters obtain most of their membership in the noise cluster and do not distort the others \cite{dave1991noise}. Third, a validity index is needed to compare solutions with different numbers of clusters. The index of \citeA{xie1991validity} relates the compactness of the clusters to the separation of their prototypes and is used by \citeA{durso2022robust}.

\section{Clustering football players}
\label{sec:lit-football}

\paragraph{Spatial and tactical clustering.} \citeA{behravan2019finding} cluster the locations from which players are active, using a particle swarm optimisation algorithm that also determines the number of clusters. Their data consist of about 93,000 objects representing the centres of the performance of players in roughly 4,900 matches in European leagues, and the method recovers positional roles. The approach uses position data only and ignores performance statistics. \citeA{narizuka2019clustering} work at the level of the team: they describe the formation of a team in a given moment as an adjacency matrix obtained from a Delaunay triangulation of the player positions and apply hierarchical clustering to these matrices, using player tracking data from Japanese league matches. Their method makes it possible to follow how a team's formation changes within and across matches, but it requires tracking data, which are not available for most leagues and seasons, and it does not profile individual players.

\paragraph{Mixed-type clustering of performance profiles.} Player performance data contain counts, rates and categorical information, so the choice of dissimilarity is central. \citeA{akhanli2023clustering} analyse the 2014--15 season of eight major European leagues (England, Spain, Italy, Germany, France, Russia, the Netherlands and Turkey). Of 3,003 players who appeared at least once, goalkeepers were excluded and 1,501 outfield players with sufficient playing time were analysed. They combine group-wise dissimilarities in which each group of variables is standardised by the standard deviation of its dissimilarities and weighted proportionally to the number of variables in the group, and they select among clustering methods and numbers of clusters with a composite index that aggregates several calibrated validity criteria. A survey among 13 football experts informs the weighting of the criteria for a second, much finer clustering with 100 to 150 clusters, intended for scouting; a coarse solution with five clusters serves broad classification.

\citeA{durso2022robust} propose the model on which this thesis is built, a robust fuzzy $C$-medoids clustering model for mixed data with a noise cluster (FCMd-MD-NC). They group the variables into types, use an appropriate dissimilarity for each type, and estimate the weight of each type during the optimisation. Applied to the 2018/19 Serie A season, with data from WhoScored, they retain 397 of 544 players, namely those with at least 200 minutes played, because per-90-minute statistics based on fewer minutes are unreliable. The best solution has three clusters and a noise cluster that contains 49 players, with $m=1.3$. They show that fuzzy memberships reveal intermediate players between the archetypes. They note that profiles of this kind are relevant for clubs in team selection and in determining the value of a player in the transfer window, and their concluding remarks point to extensions that include the temporal dimension of the playing variables and financial variables of the player and the team.

\paragraph{Longitudinal work.} The study closest to a longitudinal analysis is that of \citeA{pacifico2025fuzzy}. The author assembles, by manually reconciling statistics from Transfermarkt, FBref and WhoScored, a panel of 27 representative players who featured in European competitions and for their national teams, over the twelve seasons from 2012/13 to 2023/24, with 35 numerical variables; the data are described as a synthetic dataset. A fuzzy clustering of offensive indicators yields two profiles of mostly forwards, and a gradient boosting model with interpretability tools is used to predict whether a player exceeds a goals-per-shot threshold. The paper suggests that updating the data each season would allow analysts to monitor players' cluster transitions, but it does not analyse such transitions and does not relate them to market valuations. Salary enters only as one predictor among others.

\section{The longitudinal gap and alternative approaches}
\label{sec:lit-gap}

The studies above establish fuzzy clustering of mixed-type performance data as a rigorous way of profiling players, but with the exception of the small panel just described, each analyses a single season. A single season delivers a snapshot, whereas value is formed over time: clubs pay for expected future contributions, and the trajectory of a player may therefore be informative in a way that the current profile is not. At the same time, the valuation literature in Section~\ref{sec:lit-valuation} treats the profile through a list of individual statistics. The combination of the two, namely the movement of a player through a space of performance archetypes and its association with changes in market value, has not been studied. This is the gap addressed by the thesis, and it coincides with the extensions suggested by \citeA{durso2022robust}.

Several alternative strategies could in principle address the gap, and it is worth stating why the approach adopted here is preferred. \emph{Time-series clustering} treats each player as one multivariate time series and clusters the whole trajectories; fuzzy versions allow a series to switch between clusters \cite{durso2009autocorrelation}. This answers a different question, namely which players follow similar paths, and it assigns each player a single cluster for the whole period. It needs reasonably long and complete series, whereas the data here are an unbalanced panel of at most eight seasons in which players enter and leave the league. \emph{State-space or hidden Markov formulations} would model the evolution of a latent profile explicitly, but they rest on parametric assumptions about the dynamics that are difficult to justify with annual observations. \emph{Dimensionality reduction}, for example principal component analysis, would summarise the performance statistics in a few scores whose changes could be related to value. The scores are linear combinations of the same statistics, they do not produce interpretable archetypes, and they do not allow for partial membership; the thesis nevertheless uses them as an empirical benchmark (Chapter~\ref{ch:robustness}). The approach adopted here clusters each season separately, aligns the archetypes across seasons, and takes differences of the membership vectors. Its price is that it ignores temporal dependence in the estimation of the profiles; its advantages are transparency, direct comparability with the cross-sectional literature, and an interpretation in terms of players and archetypes that practitioners can follow.

\section{Research questions}
\label{sec:lit-questions}

The thesis pursues one main research question and two supporting questions. The main question is:

\begin{quote}
\emph{Do longitudinal transitions in fuzzy cluster membership, derived from a robust mixed-data clustering model applied to Primeira Liga player performance statistics across consecutive seasons, have incremental explanatory power for changes in transfer market valuations beyond what individual performance metrics and basic player and club characteristics capture?}
\end{quote}

The objective is explanatory and associational. The design is observational, no causal identification strategy is introduced, and findings are reported as associations. Explanatory power is assessed in sample, with tests of the added regressors and the gain in adjusted $R^2$, and out of sample, with the gain in cross-validated $R^2$ as a diagnostic of predictive relevance. The two supporting questions are descriptive and methodological:

\begin{enumerate}
	\item[RQ1.] Which player archetypes does a robust mixed-data fuzzy clustering identify in the Primeira Liga, how stable are they across seasons, and how do players move between them?
	\item[RQ2.] How sensitive is the answer to the main question to the choices that have to be made when constructing the transitions, in particular the number of clusters, the fuzziness of the memberships, the treatment of rate variables, the noise cluster, the choice of baseline, and the use of estimated rather than observed memberships in the regression?
\end{enumerate}
